\section{A polynomial kernel for the Stieltjes derivative tower}
\label{zero:section}

The rank calculation describes relations at rational arguments.  A different
question is how the functions themselves change sign as the positive real
argument varies.  Here the derivative order $k$ provides a concentration
parameter.  We keep the Laurent index $n$ fixed and let $k$ grow; this is
different from the frequently studied asymptotics of $\gamma_n(a)$ as
$n\to\infty$.

Define monic polynomials $P_n$ by
\begin{equation}\label{zero:appell}
 \frac{e^{xz}}{\Gamma(1+z)}
   =\sum_{n=0}^{\infty}P_n(x)\frac{z^n}{n!}.
\end{equation}
The generating function is entire in $z$.  With $y=x+\gamma$, its first
members are
\begin{align*}
 P_0(x)&=1, & P_1(x)&=y,\\
 P_2(x)&=y^2-\zeta(2),&
 P_3(x)&=y^3-3\zeta(2)y+2\zeta(3),\\
 P_4(x)&=y^4-6\zeta(2)y^2+8\zeta(3)y
                       +3\zeta(2)^2-6\zeta(4).
\end{align*}
They satisfy $P_n'=nP_{n-1}$ and the constructive recurrence
\begin{equation}\label{zero:recurrence}
 P_{n+1}(x)=(x+\gamma)P_n(x)
       +\sum_{r=1}^{n}(-1)^r\frac{n!}{(n-r)!}\zeta(r+1)P_{n-r}(x).
\end{equation}
This follows by differentiating \eqref{zero:appell} in $z$ and using the
Taylor series of $\log\Gamma(1+z)$ near zero.  These are the Appell
polynomials associated with the reciprocal gamma function; the use of
real-rooted Appell polynomials belongs to classical Laguerre--P\'olya
theory.  We supply the needed real-root and multiplicity argument below,
so no classification theorem is required.

\begin{theorem}[Exact Laplace kernel]\label{zero:kernel}
For $n\geq0$, $k\geq1$, and $a>0$, set
$F_{n,k}(a)=(-1)^{n+k}\gamma_n^{(k)}(a)$.  Then
\begin{equation}\label{zero:laplace}
 F_{n,k}(a)=\int_0^\infty
       \frac{t^k e^{-at}}{1-e^{-t}}P_n(\log t)\,dt.
\end{equation}
The integral and all its derivatives in $a$ converge locally uniformly
on $a>0$, and $F_{n,k}'=-F_{n,k+1}$.
\end{theorem}
\begin{proof}
Start from the classical Mellin integral for Hurwitz zeta
\cite{DLMF25}.  Differentiating $k$ times in the parameter gives, initially
for $\Re s>1$,
\[
 \partial_a^k\zeta(s,a)
   =\frac{(-1)^k}{\Gamma(s)}
      \int_0^\infty\frac{t^{s+k-1}e^{-at}}{1-e^{-t}}\,dt.
\]
For $k\geq1$ the right side is holomorphic in a neighborhood of $s=1$:
near $t=0$ its integrand is $O(t^{k-1+\Re(s-1)})$, and infinity is
controlled by $e^{-at}$.  The parameter derivative of the pole
$1/(s-1)$ is zero.  Put $s=1+z$, insert \eqref{zero:appell}, and
compare the coefficient of $z^n$ with
$(-1)^n\gamma_n^{(k)}(a)/n!$.  The bounds just stated, with finitely
many logarithmic powers or extra factors of $t$, justify every
coefficient extraction and parameter differentiation.
\end{proof}

This representation is compatible with Coffey's derivative formulas
\cite{Coffey2009} and the harmonic-number master identity in the draft.
Its useful feature here is that $P_n$ is independent of $k$: the entire
dependence on derivative order is the positive multiplier $t^k$.

\subsection{Why all kernel roots are real and simple}

\begin{lemma}\label{zero:root-preserver}
Let $p$ be a monic real polynomial with only real roots, and let $c>0$.
Then $p+cp'$ has only real roots.  A root of $p$ of multiplicity $m$
has multiplicity $m-1$ in $p+cp'$ if $m>1$ and is absent if $m=1$.
Every remaining root is simple.  Applying $n-1$ operators of the form
$1+cD$ to any real-rooted degree-$n$ polynomial therefore gives a
polynomial with $n$ distinct real roots.
\end{lemma}
\begin{proof}
If the distinct roots are $r_1<\cdots<r_\ell$, with multiplicities
$m_i$, then away from them
\[
 \frac{p'(x)}{p(x)}=\sum_{i=1}^{\ell}\frac{m_i}{x-r_i},
 \qquad
 \left(\frac{p'}p\right)'=-\sum_{i=1}^{\ell}
                                  \frac{m_i}{(x-r_i)^2}<0.
\]
The equation $p'/p=-1/c$ has exactly one simple solution in each
$(r_i,r_{i+1})$ and one in $(-\infty,r_1)$, and none in
$(r_\ell,\infty)$.  At each old root the multiplicity decreases
by one, as direct factorization shows.  This accounts for all $n$
roots.  Iteration reduces the largest possible multiplicity to one.
\end{proof}

\begin{theorem}[Simple kernel roots]\label{zero:roots}
For every $n\geq1$, $P_n$ has exactly $n$ distinct real roots.
The roots of $P_{n-1}$ strictly interlace those of $P_n$.
\end{theorem}
\begin{proof}
The Weierstrass product \cite{DLMF5} is
\[
 \frac1{\Gamma(1+z)}
   =e^{\gamma z}\prod_{j=1}^{\infty}(1+z/j)e^{-z/j}.
\]
Its finite products converge uniformly on compact sets.  At degree $n$
the corresponding Appell polynomial is
\[
 \prod_{j=1}^{M}(1+D/j)(x+\gamma-H_M)^n.
\]
Lemma~\ref{zero:root-preserver} makes it real-rooted, and convergence
of coefficients shows that $P_n$ is real-rooted and monic.  This limit
argument alone does not prove simplicity, since roots could coalesce.

To exclude coalescence, factor the first $n-1$ linear factors out of
the entire function, leaving
\[
 g_n(z)=e^{(\gamma-H_{n-1})z}
             \prod_{j=n}^{\infty}(1+z/j)e^{-z/j}.
\]
Its degree-$n$ Appell polynomial $Q_n(x)=n![z^n]e^{xz}g_n(z)$ is
again a coefficient limit of monic real-rooted polynomials.  Hence it
is real-rooted.  But now the exact finite identity
\[
 P_n=\prod_{j=1}^{n-1}(1+D/j)Q_n
\]
and Lemma~\ref{zero:root-preserver} give simplicity after taking
the limit.  The case $n=1$ is immediate.  Finally
$P_n'=nP_{n-1}$ and Rolle's theorem give strict interlacing.
\end{proof}

\section{A global bound on positive real zeros}
\label{zero:count-section}

We use a short form of the variation-diminishing principle for Laplace
transforms, proved here with Rolle's theorem.  A sign change means a
change between strict positive and negative signs; isolated zeros
without a sign change are not counted as sign changes.

\begin{lemma}[Laplace sign-variation bound]\label{zero:variation}
Let $f$ be a nonzero continuous real-valued function on $(0,\infty)$ with exactly
$m$ sign changes and with
$\int_0^\infty e^{-at}(1+t^r)|f(t)|\,dt<\infty$ for every $a>0$
and nonnegative integer $r$.  Assume $f$ is nonzero on a
set of positive measure.  Its Laplace transform
$L_f(a)=\int_0^\infty e^{-at}f(t)\,dt$ has at most $m$ zeros on
$(0,\infty)$, counted with analytic multiplicity.
\end{lemma}
\begin{proof}
The integrability hypotheses give a real-analytic Laplace transform
and justify differentiation under its integral.
For $m=0$ the transform has a fixed strict sign.  For $m>0$ choose
one sign-change point $t_0$.  Multiplication by $t_0-t$ removes
that sign change and creates no new one, so
$g(t)=(t_0-t)f(t)$ has $m-1$ sign changes.  Moreover,
\[
 \frac{d}{da}\bigl(e^{at_0}L_f(a)\bigr)=e^{at_0}L_g(a).
\]
If $L_f$ has $N$ zeros on a compact interval, counted with
multiplicity, the derivative on the left has at least $N-1$ zeros
there.  By induction $N-1\leq m-1$.  Applying this to any finite
collection of zeros proves the result on the entire half-line.
In particular, none of the transforms encountered in the induction
can vanish identically.
\end{proof}

\begin{theorem}[Uniform zero-count bound]\label{zero:at-most}
For every $n\geq0$ and every $k\geq1$, the real function
$a\mapsto\gamma_n^{(k)}(a)$ has at most $n$ positive real zeros,
counted with multiplicity.
\end{theorem}
\begin{proof}
The kernel in \eqref{zero:laplace}, before the factor $e^{-at}$,
is $t^kP_n(\log t)/(1-e^{-t})$.  Its positive factors do not change
sign, and Theorem~\ref{zero:roots} gives exactly $n$ sign changes.
All exponential moments exist by the estimates in the proof of
Theorem~\ref{zero:kernel}.  Apply Lemma~\ref{zero:variation}.
\end{proof}

The bound is independent of $k$, although the zero locations move to
infinity as $k$ increases.  We will prove that it is attained for
every sufficiently large $k$ at each fixed $n$.  There is no claim
here that it is attained at every small derivative order for $n>1$.

For reference, elementary endpoint estimates are
\begin{align}
 F_{n,k}(a)&\sim \Gamma(k+1)a^{-k-1}\bigl(\log(1/a)\bigr)^n,
                  &&a\downarrow0,\label{zero:left-end}\\
 F_{n,k}(a)&\sim (-1)^n\Gamma(k)a^{-k}(\log a)^n,
                  &&a\to\infty.\label{zero:right-end}
\end{align}
For \eqref{zero:left-end}, substitute $u=at$ and use
$(1-e^{-u/a})^{-1}\to1$ away from zero; the small-$u$ part is
controlled by $(1-e^{-v})^{-1}\leq1+1/v$.
For \eqref{zero:right-end}, the same substitution uses
$(1-e^{-u/a})^{-1}\sim a/u$ and the leading monic term of $P_n$.
Dominated convergence after division by the displayed logarithmic
power proves both statements; integrals of $u^{k-1}|\log u|^r e^{-u}$
are finite.  The interpretations for $n=0$ are immediate.

\section{Uniform expansions and the motion of every zero}
\label{zero:asymptotic-section}

Put $K=k+1$ and
\[
 h_n(t)=\frac{P_n(\log t)}{1-e^{-t}},\qquad
 \mathcal H_{n,K}(t)=\frac{(K/t)^K}{\Gamma(K)}F_{n,K-1}(K/t).
\]
If $U_K$ has gamma density
$K^K u^{K-1}e^{-Ku}/\Gamma(K)$ on $u>0$, then
\begin{equation}\label{zero:expectation}
 \mathcal H_{n,K}(t)=\mathbb E\,h_n(tU_K),
 \qquad \mathbb E U_K=1.
\end{equation}
This is an exact identity.  It identifies the useful scale $a=K/t$
and turns the asymptotic calculation into an expansion around $U_K=1$.

\begin{theorem}[All-order local expansion]\label{zero:all-orders}
Fix $n\geq0$, $M\geq0$, a compact interval $I\subset(0,\infty)$,
and a nonnegative integer $r$.  Uniformly for $t\in I$, together with
the first $r$ derivatives in $t$,
\begin{equation}\label{zero:expansion}
 \mathcal H_{n,K}(t)=\sum_{j=0}^{M}\frac{\mathcal D_jh_n(t)}{K^j}
                         +O_{n,M,r,I}(K^{-M-1}).
\end{equation}
Here $\mathcal D_0h=h$, and for $j\geq1$ the finite formula is
\begin{equation}\label{zero:D-general}
 \mathcal D_jh(t)=
 \sum_{\substack{m_2,m_3,\ldots\geq0\\
                 \sum_{\ell\geq2}(\ell-1)m_\ell=j}}
 \frac{t^R h^{(R)}(t)}{
             \prod_{\ell\geq2}m_\ell!\,\ell^{m_\ell}},
 \qquad R=\sum_{\ell\geq2}\ell m_\ell.
\end{equation}
In particular,
\begin{align}
 \mathcal D_1h&=\tfrac12t^2h'',\label{zero:D1}\\
 \mathcal D_2h&=\tfrac13t^3h'''+\tfrac18t^4h^{(4)},\label{zero:D2}\\
 \mathcal D_3h&=\tfrac14t^4h^{(4)}+\tfrac16t^5h^{(5)}
                                      +\tfrac1{48}t^6h^{(6)}.
\end{align}
The constants in the remainder can be bounded using suprema of finitely
many derivatives of $h_n$ on an enlarged compact interval and explicit
gamma moments and tail bounds.  No uniformity in an increasing $n$
is asserted.
\end{theorem}
\begin{proof}
The cumulants of $U_K-1$ are zero in degree one and
$(\ell-1)!K^{1-\ell}$ in degree $\ell\geq2$.  Expanding its moment
generating function therefore gives the exact central-moment identity
\begin{equation}\label{zero:moments}
 \frac{\mathbb E(U_K-1)^R}{R!}
 =\sum_{\substack{\sum_{\ell\geq2}\ell m_\ell=R}}
       \prod_{\ell\geq2}
       \frac{K^{-(\ell-1)m_\ell}}{m_\ell!\,\ell^{m_\ell}}.
\end{equation}
For $R=0$ the value is one, and for $R=1$ it is zero.
When $R\geq2$, the smallest possible power of $K^{-1}$ is
$\lceil R/2\rceil$.

Expand $h_n(tu)$ about $u=1$ to degree $2M+1$, on
$1/2\leq u\leq2$.  Its remainder is bounded uniformly by a constant
times $|u-1|^{2M+2}$.  The expectation of this even power is
$O(K^{-M-1})$ by \eqref{zero:moments}.  The coefficient of $K^{-j}$
in the expected Taylor polynomial is precisely
\eqref{zero:D-general}; terms with $j>M$ are $O(K^{-M-1})$.

For completeness, the discarded tails are exponentially small, even
though $h_n$ is unbounded at zero.  Chernoff's bound for the gamma
moment generating function gives, for fixed $v\ne1$ on its relevant
side of one,
\[
 \Pr(U_K\geq v)\ \text{or}\ \Pr(U_K\leq v)
       \leq \exp\{-K(v-1-\log v)\}.
\]
For each fixed derivative order, $u^r h_n^{(r)}(tu)$ is bounded
uniformly on $t\in I$ by $C(u^{-B}+u^B)$ for some fixed $B$.
The gamma moments
$\mathbb E U_K^v=K^{-v}\Gamma(K+v)/\Gamma(K)$ are uniformly bounded
for fixed positive or negative $v$ once $K$ is sufficiently large.
Cauchy--Schwarz therefore bounds the tail expectations by
$Ce^{-cK}$.  The same reasoning replaces truncated central moments
by their full values.  Finally differentiate
\eqref{zero:expectation} in $t$, obtaining
$\mathbb E[U_K^r h_n^{(r)}(tU_K)]$, and repeat the argument.
This proves the stated $C^r$ uniformity and the expansion.
\end{proof}

\begin{theorem}[Saturation and asymptotic locations]\label{zero:locations}
Fix $n\geq1$, and let $\rho$ run over the distinct real roots of $P_n$.
For all sufficiently large $k$, $\gamma_n^{(k)}$ has exactly $n$
positive real zeros, all simple.  One zero is associated with each
$\rho$.  With $t_\rho=e^\rho$, it has the expansion
\begin{equation}\label{zero:location-first}
 a_{\rho,k}=\frac{k+1}{t_\rho}
       +\frac{1}{2t_\rho}
           \left(\frac{P_n''(\rho)}{P_n'(\rho)}-1\right)
       -\frac1{e^{t_\rho}-1}+O_n(k^{-1}).
\end{equation}
Each zero has a full expansion in descending integer powers of $k+1$.
In particular $a_{\rho,k}/k\to e^{-\rho}$.
\end{theorem}
\begin{proof}
The zeros of $h_n$ are the simple points $t_\rho$.  By the uniform
$C^1$ version of Theorem~\ref{zero:all-orders}, each has a disjoint
neighborhood containing a unique simple zero $t_{\rho,K}$ of
$\mathcal H_{n,K}$ for all sufficiently large $K$.  Thus there are
at least $n$ positive zeros.  Theorem~\ref{zero:at-most} excludes
every additional zero, including any outside the compact intervals
used in this local argument.

Write $t=t_\rho$, $h=h_n$, and
$A(t)=t^2h''(t)/2$.  Expansion of the zero equation gives
\[
 t_{\rho,K}=t-\frac{A(t)}{K h'(t)}+O(K^{-2}),
 \qquad
 a_{\rho,k}=\frac K{t_{\rho,K}}
              =\frac Kt+\frac{h''(t)}{2h'(t)}+O(K^{-1}).
\]
At a root of $P_n(\log t)$, direct differentiation gives
\[
 \frac{h''(t)}{h'(t)}
  =\frac1t\left(\frac{P_n''(\rho)}{P_n'(\rho)}-1\right)
                           -\frac2{e^t-1},
\]
which proves \eqref{zero:location-first}.  At every subsequent order
the new coefficient of $t_{\rho,K}$ is multiplied by the same nonzero
number $h'(t)$; the all-order expansion and a Taylor remainder therefore
determine it recursively.  Applying the mean value theorem to the
residual of a truncated zero expansion proves the stated remainder at
each finite order.
\end{proof}

The distinction between local concentration and global zero counting is
essential: concentration alone finds $n$ nearby zeros but cannot rule
out additional zeros on scales escaping the chosen compact interval.
The variation bound supplies precisely that missing global step.

\paragraph{An explicit next correction.}
Let $b_j=h_n^{(j)}(t_\rho)/h_n'(t_\rho)$ for $j=2,3,4$ and set
\begin{equation}\label{zero:next-coefficient}
 c_\rho=t_\rho\left(\frac{b_3}{3}-\frac{b_2^2}{4}\right)
      +\frac{t_\rho^2}{8}
                    \left(b_4-2b_2b_3+b_2^3\right).
\end{equation}
Then the more accurate formula is
\begin{equation}\label{zero:location-second}
 a_{\rho,k}=\frac{k+1}{t_\rho}+\frac{b_2}{2}
                         +\frac{c_\rho}{k+1}+O_n(k^{-2}).
\end{equation}
Indeed, with $B=\mathcal D_2h$, the first two coefficients of
$t_{\rho,K}=t+u_1/K+u_2/K^2+\cdots$ satisfy
$u_1=-A/h'$ and
$u_2=-(h''u_1^2/2+A'u_1+B)/h'$.  Expanding $K/t_{\rho,K}$
and collecting terms gives \eqref{zero:next-coefficient}.

\begin{corollary}[Persistence and interlacing after saturation]
\label{zero:k-interlace}
Fix $n\geq1$. If $\gamma_n^{(k)}$ has $n$ distinct positive
zeros at some $k\geq1$, then the same is true at order $k+1$.
Writing the zeros in increasing order, one has
\[
 a_{j,k}<a_{j,k+1}<a_{j+1,k}\quad(1\leq j<n),
 \qquad a_{n,k}<a_{n,k+1}.
\]
In particular, these conclusions hold at every sufficiently large $k$.
\end{corollary}
\begin{proof}
Rolle's theorem and $F_{n,k}'=-F_{n,k+1}$ give a zero between each
consecutive pair of zeros of $F_{n,k}$.  Beyond its last zero,
$F_{n,k}$ has a fixed nonzero sign and tends to zero by
\eqref{zero:right-end}; it therefore has an interior extremum there.
This gives $n$ distinct zeros of $F_{n,k+1}$.  Its zero bound
excludes all others, including coincidences with a simple zero of
$F_{n,k}$, and establishes the claimed ordering.
\end{proof}

\subsection{The first Laurent index: a global half-unit bracket}

For $n=1$, the kernel has its sole sign change at $t_0=e^{-\gamma}$.
Here the conclusions can be made global in the derivative order.
Coffey already proved that $\gamma_1$ has a unique maximum at
$a\approx1.39112$ \cite[Proposition 5]{Coffey2009}; this is the
$k=1$ uniqueness case below.  The present statement extends the
description to every positive derivative order and supplies the
uniform bracket and moving-zero expansion.

\begin{theorem}[Unique zero and an explicit bracket]\label{zero:first}
For every integer $k\geq1$, $\gamma_1^{(k)}(a)$ has a unique positive
zero $a_k$, which is simple.  With $H_0=0$,
\begin{equation}\label{zero:half-unit}
 e^{H_{k-1}}<a_k<e^{H_{k-1}}+\tfrac12.
\end{equation}
The zeros are strictly increasing in $k$.  Their asymptotic expansion
starts with
\begin{equation}\label{zero:first-asymptotic}
 a_k=e^\gamma\left(k+\tfrac12\right)
                     -\frac1{e^{e^{-\gamma}}-1}+O(k^{-1}).
\end{equation}
Formula~\eqref{zero:location-second} supplies the next coefficient
and an $O(k^{-2})$ remainder.
\end{theorem}
\begin{proof}
The endpoint signs in \eqref{zero:left-end}--\eqref{zero:right-end}
and the bound of one zero prove existence, uniqueness, and simplicity.
To obtain the bracket, normalize the positive density
$t^k e^{-at}/(1-e^{-t})$.  The sign of $F_{1,k}(a)$ is the sign of
$\mathbb E(\log T+\gamma)$ under this density.

Relative to a gamma density of shape $k$ and rate $a$, the tilting
factor is
$r(t)=t/(1-e^{-t})$, which is strictly increasing: its derivative has
numerator $1-e^{-t}(1+t)>0$.  A strictly increasing tilt strictly
increases the expectation of the strictly increasing function
$\log t$.  This elementary fact follows from
\[
 2\operatorname{Cov}(u(X),v(X))
   =\mathbb E[(u(X)-u(Y))(v(X)-v(Y))]>0
\]
for independent $X,Y$ with the same continuous distribution and
strictly increasing $u,v$.  Consequently
\[
 \mathbb E\log T>\psi(k)-\log a
                 =H_{k-1}-\gamma-\log a.
\]
At the zero this gives $a_k>e^{H_{k-1}}\geq1$.

For $a>1/2$, compare instead with a gamma density of shape $k$
and rate $a-1/2$.  The tilting factor is
$s(t)=t/(2\sinh(t/2))$, which is strictly decreasing because
$(t/2)\cosh(t/2)>\sinh(t/2)$.  The covariance inequality reverses:
\[
 \mathbb E\log T<\psi(k)-\log(a-1/2).
\]
Applying it at $a_k>1$ gives $a_k<e^{H_{k-1}}+1/2$.

Finally, at a zero the signed kernel $f(t)$ in
\eqref{zero:laplace} satisfies $\int e^{-a_kt}f(t)\,dt=0$ and
$(t-t_0)f(t)>0$ except at $t_0$.  Thus
\[
 F_{1,k+1}(a_k)=\int_0^\infty (t-t_0)e^{-a_kt}f(t)\,dt>0.
\]
Since $F_{1,k+1}$ crosses from positive to negative at its unique
zero, $a_{k+1}>a_k$.  The root of $P_1$ is $-\gamma$, so
\eqref{zero:first-asymptotic} is Theorem~\ref{zero:locations}.
\end{proof}

The interval in \eqref{zero:half-unit} is an analytic certificate for
bracketing a unique root, valid even at $k=1$.  Its width is uniformly
$1/2$; no assertion that this is the best possible width is made.
The accompanying computation uses this bracket rather than guessing
an initial zero from a graph.

\begin{figure}[tbp]
\centering
\includegraphics[width=0.98\linewidth]{figures/stieltjes_profiles.pdf}
\caption{Normalized derivative profiles
$(1-e^{-e^x})\mathcal H_{n,K}(e^x)$ approach $P_n(x)$ as $K=k+1$
increases.  The plots illustrate $n=1,2,3$.  Quadrature produces the
curves; Theorems~\ref{zero:at-most} and \ref{zero:locations}, rather
than the plots, prove the zero counts and asymptotic completeness.}
\label{zero:profile-figure}
\end{figure}
